

\loesung{

The implementations of the solution in R or in Python can be found on Moodle in the files "\textit{CART\_sol.R}" and "\textit{CART\_sol.py}".
As described in the exercise, the main idea is not to compute the sums for the two potential child nodes in every step.
In the following, we present two ways to work around this.

\textbf{Solution 1 (Cumulative sums):}
The first way is to compute the cumulative sums of the target vector, that is, compute a vector which at position $k$ contains the entry $\sum_{i=1}^k y_i$ (or $y_i^2$ respectively).
Then, for each split point, find the corresponding index of the feature of interest, and look up the respective sums in the precomputed vector in $O(1)$ time.

\begin{algorithm}[H]
\caption{Efficient split search -- cumulative sums (complexity: $O(n)$)}
\begin{algorithmic}[1]
\Require sorted targets $y_{(1)}, \ldots, y_{(n)}$ and candidate split values $S$
\State $\texttt{cumsum\_y}[k] \gets \sum_{i=1}^{k} y_{(i)}$ for $k = 1, \ldots, n$ \Comment{Precompute once in $O(n)$}
\State $\texttt{cumsum\_y2}[k] \gets \sum_{i=1}^{k} y_{(i)}^2$ for $k = 1, \ldots, n$
\For{each split value $s \in S$}
    \State Find index $k$ such that $x_{(k)} \le s < x_{(k+1)}$
    \State $\text{sum\_L} \gets \texttt{cumsum\_y}[k]$, \quad $\text{sum\_R} \gets \texttt{cumsum\_y}[n] - \text{sum\_L}$ \Comment{$O(1)$ update}
    \State $\text{sum\_L2} \gets \texttt{cumsum\_y2}[k]$, \quad $\text{sum\_R2} \gets \texttt{cumsum\_y2}[n] - \text{sum\_L2}$
    \State Compute $\text{Var}_L = \text{sum\_L2}/n_L - (\text{sum\_L}/n_L)^2$, analogously $\text{Var}_R$
    \State Compute loss reduction $\Delta$
\EndFor
\State \textbf{return} split value $s^*$ with largest $\Delta$
\end{algorithmic}
\end{algorithm}
This solution is implemented as ``Solution 1'' in the code.

\textbf{Solution 2 (Running sums):}
The other, more direct way is to maintain two sums: one for the potential left child of the current split point and one for the potential right child.
These two sums are then updated incrementally in every step, so that the left sum always increases and the right sum always decreases.

This method is implemented as ``Solution 2'' in the code.

\begin{algorithm}[H]
\caption{Efficient split search -- running sums (complexity: $O(n)$)}
\begin{algorithmic}[1]
\Require sorted targets $y_{(1)}, \ldots, y_{(n)}$ and candidate split values $S$ (sorted)
\State $\text{sum\_L} \gets 0$, \quad $\text{sum\_R} \gets \sum_{i=1}^{n} y_{(i)}$ \Comment{Precompute once in $O(n)$}
\State $\text{sum\_L2} \gets 0$, \quad $\text{sum\_R2} \gets \sum_{i=1}^{n} y_{(i)}^2$
\State $n_L \gets 0$, \quad $n_R \gets n$
\For{each split value $s$ (in sorted order)}
    \State Move observations from right to left: $\text{sum\_L} \mathrel{+}= y_{(k)}$, $\text{sum\_R} \mathrel{-}= y_{(k)}$ \Comment{$O(1)$ update}
    \State Update $\text{sum\_L2}$, $\text{sum\_R2}$, $n_L$, $n_R$ analogously
    \State Compute $\text{Var}_L$, $\text{Var}_R$, and loss reduction $\Delta$
\EndFor
\State \textbf{return} split value $s^*$ with largest $\Delta$
\end{algorithmic}
\end{algorithm}

Both these solutions, in principle, achieve the desired result, namely that the whole algorithm of finding the optimal split point has a computational complexity of $O(dn)$ instead of $O(d n^2)$.
Nevertheless, the two implementations differ in another point, which may be practically even more relevant.
Namely, the ``Solution 1'' relies on vectorized operations, acting in parallel on the entire vector of observations or potential split points. These vectorized operations, in theory, require an effort of order $O(n)$ each (linear in the vector's length), because they perform one operation per vector element, but they are much faster in practice due to parallelization.

The ``Solution 2'' does not use such operations at all and therefore strictly fulfills the theoretical runtime requirement, whereas ``Solution 1'' strictly speaking does not.
Nevertheless, at least for the R implementation, ``Solution 1'' actually runs faster in practice.

The solution also includes the calculation of a deeper tree that reflects the structure of the DGP used.

NOTE: This exercise also serves as preparation for the FAST algorithm used in explainable boosting machines (EBMs) introduced later in the lecture.

}